\section{Tomographic Initialization}
\label{sec:tomography}

Rendering a 3D Gaussian field produces, for every view, a field of projected 2D
Gaussians.  The captures of \cref{sec:fields} are fields of exactly this type, so each
per-view 2D Gaussian field is a projection of the same unknown 3D radiance density.  This
suggests obtaining the initialization by inverting the projection.  We therefore pose 3D
initialization as reconstruction from projections instead of discrete correspondence
search.  Gaussians keep the inversion tractable in closed form, since the back-projection
of a 2D Gaussian is an analytic elongated 3D Gaussian and fusions of Gaussians remain
Gaussian.  The resulting method, called Beam Fusion, is deterministic, RGB-free and
correspondence-free.  Association emerges from density overlap instead of discrete
matching, and no voxel grid is built.

\subsection{Back-projection}
\label{sec:beams}

Fix view $v$ with camera center $\bm{o}_v$ and a fitted component $(\bm{m}_j, C_j, a_j,
\bm{c}_j)$.  Its center pixel defines a world ray direction $\bm{d}_{vj}$.  At a
hypothesised camera depth $t$ the projective Jacobian of the pinhole map at the point
$\bm{p} = t\,\bm{r}_{vj}$ in the camera frame is
\begin{equation}
  J(\bm{p}) \;=\;
  \begin{pmatrix}
    f_x / z & 0 & -f_x\, x / z^2\\
    0 & f_y / z & -f_y\, y / z^2
  \end{pmatrix},
  \qquad \bm{p} = (x, y, z)\T .
  \label{eq:jacobian}
\end{equation}
Let $B = (\bm{b}_1, \bm{b}_2, \bm{r})$ be an orthonormal camera-frame basis whose third
column is the ray direction.  The fitted 2D covariance is lifted to the ray-orthogonal
tangent plane through the inverse of the projected tangent Jacobian as
\begin{equation}
  \Sigma_{\perp}
  \;=\;
  \big(J B_{:,1:2}\big)^{-1}\, C_j\, \big(J B_{:,1:2}\big)^{-\mathsf{T}}.
  \label{eq:tangent-lift}
\end{equation}
This lift is exact for the center ray.  It answers which tangent-plane covariance would
have projected to the observed $C_j$.  Along the ray the depth of the component is
unknown within the working depth range $[t_{\mathrm{lo}}, t_{\mathrm{hi}}]$, which we
obtain by intersecting the ray with the scene bounding box, clamped to a near plane.  The
beam assigns the ray direction a long variance $L^2$ with half-length
$L = \tfrac12 (t_{\mathrm{hi}} - t_{\mathrm{lo}})$.  The beam is the 3D Gaussian with
world-frame precision
\begin{equation}
  \Lambda_{vj}(t)
  \;=\;
  R_v\T \, B \begin{pmatrix} \Sigma_{\perp}^{-1} & 0 \\ 0 & 1/L^2 \end{pmatrix} B\T R_v .
  \label{eq:beam-precision}
\end{equation}
It is certain transversally, exactly as the image measured, and agnostic along the ray.

\subsection{Pair seeding}
\label{sec:seeding}

For components $i$ in view $u$ and $j$ in view $v$, the closest points between their
center rays have closed-form depths $(s, t)$ given by the standard two-ray least-squares
solution.  Parallel rays are rejected.  A candidate 3D location is the midpoint of the
closest-point segment.  Candidates pass three analytic tests,
\begin{equation}
  \underbrace{\frac{\norm{\bm{p}_u(s) - \bm{p}_v(t)}^2}
       {\sigma_u^2(s) + \sigma_v^2(t)} \le 3^2}_{\text{transverse gate}},
  \qquad
  \underbrace{\norm{\bm{c}_i - \bm{c}_j} \le 0.35}_{\text{color gate}},
  \qquad
  \underbrace{s \in [t^u_{\mathrm{lo}}, t^u_{\mathrm{hi}}],\;
              t \in [t^v_{\mathrm{lo}}, t^v_{\mathrm{hi}}]}_{\text{depth validity}},
  \label{eq:gates}
\end{equation}
where $\sigma_u(s)$ is the mean pixel footprint of component $i$ scaled to world units at
depth $s$.  The ray separation is thus measured against both beams' transverse
footprints, and large splats tolerate larger separations.  Each surviving seed receives
the first-order Gaussian product weight
\begin{equation}
  w \;=\; \tfrac12 (a_i + a_j)\,
          \exp\!\big({-\tfrac12\, g}\big)\,
          \exp\!\big({-\norm{\bm{c}_i - \bm{c}_j}^2 / (2\sigma_c^2)}\big),
  \qquad \sigma_c = 0.25,
  \label{eq:seed-weight}
\end{equation}
with $g$ the transverse gate value, so negligible cross terms are never materialized.
View pairs are processed in increasing camera-distance order.  Our production path
streams all pairs through a bounded 3D voxel table which retains only the strongest seed
per voxel with deterministic tie-breaking.  The seed storage is thus capped independently
of the scene size.

\subsection{Covariance intersection}
\label{sec:ci-fusion}

Contributing beams of one physical splat are correlated observations, since they look at
the same surface through overlapping information.  The obvious fusion rule, the Gaussian
product $\Lambda = \sum_k \Lambda_k$, cannot be applied here.  It double counts every axis
observed by more than one view.  For example, two orthogonal views of an isotropic blob
would fuse to half its variance.  We use equal-weight covariance
intersection~\citep{julier1997_ci} instead,
\begin{equation}
  \Lambda \;=\; \frac{1}{K} \sum_{k=1}^{K} \Lambda_k,
  \qquad
  \bm{m} \;=\; \Lambda^{-1} \Big( \frac{1}{K} \sum_{k=1}^{K} \Lambda_k\, \bm{m}_k \Big),
  \label{eq:ci}
\end{equation}
where $\bm{m}_k$ is the world-space mean of beam $k$ at its implied depth.  The common
$1/K$ leaves the fused mean identical to the product rule but keeps the covariance
conservative.  Covariance intersection is exact on directions all views share and never
overconfident elsewhere, which is the standard consistency guarantee under unknown
correlation.  The main failure risk of the construction is thus false contributor
association and not an algebraic missing factor.  Contributor counts, gates and lineage
are therefore first-class outputs, see \cref{sec:reduction}.

\subsection{Minimum view count}
\label{sec:observability}

Two views are not sufficient to recover the fused covariance.  Fix a 3D Gaussian with
triangulated mean and EWA covariance projection $C_v = A_v \Sigma A_v\T$ into views
$v = 1, 2$ with world ray directions $\bm{d}_1, \bm{d}_2$ through the mean.  Each
projection Jacobian annihilates its own viewing ray, $A_v \bm{d}_v = 0$.  The symmetric
matrix $\Delta = \bm{d}_1 \bm{d}_2\T + \bm{d}_2 \bm{d}_1\T$ therefore projects to zero in
both views, due to
$A_1 \Delta A_1\T = (A_1\bm{d}_1)(A_1\bm{d}_2)\T + (A_1\bm{d}_2)(A_1\bm{d}_1)\T = 0$ and
the symmetric identity for the second view.  A dimension count over
$\Sym(3) \cong \R^6$ shows that two generic views determine only five of the six
covariance coordinates and that $\Delta$ spans the missing direction.  A third generic
view removes this null direction.  We verified the rank statement symbolically and
numerically in our test suite.\evidence{ara/logic/claims.md C23,
tests/test_inverse_projection_fiber.py}  A two-view fusion would thus carry an arbitrary
covariance component along the shared null mode and not merely a noisy one.  Beam Fusion
therefore requires $K \ge 3$ contributing views per output component, turning the
algebraic blind spot into a hard admission rule.

\subsection{Fold-in and reduction}
\label{sec:reduction}

Each surviving seed projects its provisional 3D midpoint into every remaining view and
absorbs at most one nearest fitted component per view, subject to a $3\sigma$ pixel gate
relative to the candidate component's own footprint, the color gate of \cref{eq:gates}
and a positive-depth test.  Components with fewer than three contributors are discarded,
see \cref{sec:observability}.  Each retained component is fused over all its contributors
with \cref{eq:ci}, weighted by $w \cdot K$, and reduced in two exact steps.  First,
components with identical contributor signatures collapse to the strongest instance.
Second, a per-voxel non-maximum suppression keeps the strongest component per cell.  Both
steps are selection and not moment matching, so the fused covariances survive unblurred.
Covariance eigenvalues are clamped to world bounds $[\sigma_{\min}^2, \sigma_{\max}^2]$.
Colors initialize as amplitude-weighted contributor means.  Opacity initializes at a
constant $0.1$, since \cref{sec:refinement} shows that no opacity can be inferred from
the captures.  The output is a standard 3D Gaussian set plus complete contributor
lineage, i.e.\ a CSR map from every 3D component to its supporting 2D splats and implied
depths, together with per-view unmatched-splat lists.  Downstream refinement, diagnostics
and future densification consume this by-product.  Each run also records the
contributor-count histogram, evaluated ray pairs, gate survival counts, seed-voxel
occupancy and stage timings.\evidence{src/rtgs/lift/beam_fusion.py}

\begin{figure}[t]
  \centering
  \figslot{6.5cm}{%
    \textbf{Figure 6, geometry of the construction.}  Two panels.  \textbf{Left:} one 2D
    Gaussian in two views, back-projected as two elongated beams with the transverse
    footprint of \cref{eq:tangent-lift} and the long axis $L$ from the depth interval,
    their closest-ray midpoint, the transverse gate as shaded tolerance tubes and the
    fused 3D Gaussian at the intersection.  Contrast the conservative covariance
    intersection ellipsoid with the overconfident naive product ellipsoid.
    \textbf{Right:} the fold-in step, with the provisional point projecting into a third
    view, the $3\sigma$ pixel gate around the nearest 2D component and the resulting
    three-contributor lineage.  Draw as a TikZ schematic and compute exact ellipse shapes
    from a toy two-camera configuration.}
  \caption{\redcaption{Beam back-projection, gating and covariance-intersection fusion.}}
  \label{fig:beam-geometry}
\end{figure}

\begin{figure}[t]
  \centering
  \figslot{6cm}{%
    \textbf{Figure 7, initialization comparison.}  Renders of the initial 3D Gaussian
    sets on one scene for the tomographic initialization, the image-free placement
    control, calibration-bounded random and, if available, the COLMAP sparse
    initialization, all at the same count budget.  \textbf{Top row:} front view.
    \textbf{Bottom row:} side or top view, which exposes along-ray elongation and
    floaters that front views hide.  Color the tomographic panel by contributor count
    (3, 4, 5 and more views).  Uses existing development artifacts under \repopath{runs/}
    or a rerun of the initializers on the development scene, rendered with the repository
    viewer.}
  \caption{\redcaption{Initial 3D Gaussians per initializer, front and side views.
    Producible from checked-in initializers.}}
  \label{fig:init-visual}
\end{figure}
